\documentclass[11pt]{article}
\usepackage[T1]{fontenc}
\usepackage{lmodern}
\usepackage[margin=1in]{geometry}
\usepackage{amsmath,amssymb,amsthm}
\usepackage{microtype}
\usepackage[hidelinks]{hyperref}
\newtheorem{theorem}{Theorem}
\newtheorem{lemma}[theorem]{Lemma}
\newcommand{\R}{\mathbb R}
\newcommand{\N}{\mathbb N}
\newcommand{\Ran}{\operatorname{Ran}}
\newcommand{\rank}{\operatorname{rank}}
\newcommand{\supp}{\operatorname{supp}}
\newcommand{\dom}{\operatorname{dom}}
\newcommand{\norm}[1]{\left\lVert #1\right\rVert}
\newcommand{\Hpp}{\mathcal H_{\mathrm{pp}}}
\title{Pure point Dirichlet spectrum on quasi-conical domains}
\author{Alec Kriebel\\
  \small Independent researcher\\
  \small ORCID: \href{https://orcid.org/0009-0001-9320-500X}{0009-0001-9320-500X}}
\date{15 October 2026\\ \small Research note, version 1.0}
\begin{document}
\maketitle
\begin{abstract}
For every dimension at least two, we construct a connected open
quasi-conical domain whose Dirichlet Laplacian has a complete orthonormal
eigenbasis, although its spectrum and essential spectrum are the entire
nonnegative half-line. Its positive eigenvalues are dense, zero is not
an eigenvalue, and both continuous spectral subspaces vanish.
We use the expanding-cube and shrinking-window geometry of
Krej\v{c}\'i\v{r}\'ik and Lotoreichik, choosing both cube sizes and positive
apertures inductively. The additional construction transports a growing
family of complete finite spectral clusters with summable norm errors.
A repeated dense test family proves completeness and excludes singular
continuous spectrum. This answers the existential spectral-type question
in their discussion and the corresponding Oberwolfach contribution.
\end{abstract}

\noindent\textbf{Status and disclosure.}
AI tools were used extensively in developing, reproducing, checking and
adversarially reviewing the argument and preparing this note.
This is an unrefereed preprint. Internal AI-assisted review is not external
human peer review or proof-assistant certification.
The source and verification package records the checks and their limits.

\section{The question and the precise result}

An open set is \emph{quasi-conical} if it contains Euclidean balls of
arbitrarily large radius. On any open $\Omega\subset\R^d$, let $A_\Omega$
be the nonnegative self-adjoint Dirichlet Laplacian defined by the closed
form
\[
 q_\Omega[u]=\int_\Omega|\nabla u|^2,\qquad
 \dom q_\Omega=H_0^1(\Omega).
\]
This definition imposes no regularity assumption on the boundary.
Krej\v{c}\'i\v{r}\'ik and Lotoreichik \cite{KL} constructed connected
quasi-conical domains with an embedded eigenvalue and with no absolutely
continuous spectrum. Their Section~5 asks whether the singular continuous
part can also be removed. Lotoreichik's Oberwolfach contribution
\cite[pp.~2077--2078]{OWR} poses the same existential question; it is
repeated in \cite[Open Problem~1, PDF p.~16]{K26}.

\begin{theorem}\label{thm:main}
For every integer $d\ge2$ there is a connected open quasi-conical set
$\Omega\subset\R^d$, consisting of expanding cubes joined through positive
shrinking windows in their common faces, such that $A=A_\Omega$ has a
complete orthonormal eigenbasis and
\[
 \sigma(A)=\sigma_{\mathrm{ess}}(A)
   =\overline{\sigma_{\mathrm p}(A)}=[0,\infty),
 \qquad \ker A=\{0\},\qquad
 \mathcal H_{\mathrm{ac}}(A)=\mathcal H_{\mathrm{sc}}(A)=\{0\}.
\]
Both the cube widths and the window widths are chosen in the construction.
\end{theorem}

Pure point spectral type means completeness of the eigenvectors.
It does not assert that every point of $[0,\infty)$ is an eigenvalue.
The theorem concerns zero-thickness interface screens with open windows.
It makes no assertion for arbitrary preassigned cube sizes, smooth
boundaries, Neumann conditions or positive-length connecting passages;
no effective numerical aperture bounds are supplied.

The geometry, freedom to choose future cube sizes and finite-window
continuity below come from \cite[Sections~2.2 and~3]{KL}.
Whole spectral families already handle multiplicities in
\cite[Proposition~1.4]{HSS}; see also the explicit warning in
\cite[Remark~3.1]{KL}. The additional step here is the exhaustive geometric
transport and its completeness estimate. Neither finite-cluster
perturbation theory nor diagonalization on finite reducing subspaces is
claimed as a new abstract principle.

\section{A bounded shrinking-window step}

Write $Q_a(x)=xe_1+(-a,a)^d$. Suppose that a bounded connected partial
tower $O_n$ ends in $Q_{a_n}(x_n)$, with all its earlier windows fixed.
For $a>a_n+1$, put $x=x_n+a_n+a$, $p=(x_n+a_n)e_1$, and
\[
 D_0=O_n\cup Q_a(x),\qquad
 D_\delta=D_0\cup\bigl(p+(-\delta,\delta)^d\bigr),
 \qquad 0<\delta<a_n .
\]
The connector away from its middle face is already inside the adjacent
cubes. Thus $D_\delta\setminus D_0$ is precisely a patch of that face:
the $L^2$ spaces are canonically the same, while the form domains differ.
The full connector cube makes $D_\delta$ open and connected.

\begin{lemma}\label{lem:window}
On this common $L^2$ space,
\[
 \norm{(A_{D_\delta}+1)^{-1}-(A_{D_0}+1)^{-1}}\longrightarrow0
 \quad(\delta\downarrow0).
\]
All these bounded-domain Dirichlet operators have compact resolvent and
strictly positive eigenvalues.
\end{lemma}

\begin{proof}
This is the finite-window application of
\cite[Propositions~2.2--2.3 and Section~3]{KL}. We give a direct form
argument to make the dimension and boundary hypotheses explicit.
Choose one bounded box $B$ containing the fixed finite geometry.
Zero extension realizes $V_\delta=H_0^1(D_\delta)$ and
$V_0=H_0^1(D_0)$ as closed subspaces of $H_0^1(B)$.
They decrease with $\delta$, and
\begin{equation}\label{eq:intersection}
 \bigcap_{0<\delta<a_n}V_\delta=V_0.
\end{equation}
Indeed, for bounded $u$ in this intersection choose a smooth cutoff
$\zeta_r$ that vanishes near $p$ and is one outside a shrinking
neighborhood. Choose $\delta$ so small that the connector is inside the
zero neighborhood. Multiplying smooth compactly supported approximants
to $u$ in $D_\delta$ by $\zeta_r$ shows that $\zeta_r u\in V_0$.
For $d\ge3$ use a cutoff between radii $r$ and $2r$; its squared gradient
norm is $O(r^{d-2})$. For $d=2$, choose a smooth $\theta$ equal to zero
on $(-\infty,0]$ and one on $[1,\infty)$, and use
\[
 \zeta_r(y)=
 \theta\!\left(\frac{\log(|y-p|/r^2)}{\log(1/r)}\right)
 \quad(r^2<|y-p|<r),
\]
continued by zero inside and one outside. Its gradient energy is
\[
 \int_{\R^2}|\nabla\zeta_r|^2
 =\frac{2\pi}{\log(1/r)}
       \int_0^1|\theta'(s)|^2\,ds\longrightarrow0.
\]
In either case $\zeta_r u\to u$ in $H^1(B)$ for bounded $u$:
use dominated convergence and
$\norm{u\nabla\zeta_r}_2\le\norm{u}_\infty\norm{\nabla\zeta_r}_2$.
Truncate the real and imaginary parts first to treat arbitrary $u$;
Sobolev truncation preserves each $V_\delta$ and converges in $H^1$.
Closedness proves \eqref{eq:intersection}.

If $\delta_j\to0$ and $u_j\in V_{\delta_j}$ converges weakly in $H^1(B)$,
its limit belongs to every fixed $V_s$, hence to $V_0$.
Every $v\in V_0$ is its own recovery sequence.
Now write $R_\delta=(A_{D_\delta}+1)^{-1}$ and
$R_0=(A_{D_0}+1)^{-1}$. For bounded inputs $f_j$,
the solutions $R_{\delta_j}f_j$ are uniformly bounded in $H_0^1(B)$.
If norm convergence failed, select unit inputs with a positive lower
bound on $\norm{(R_{\delta_j}-R_0)f_j}_2$.
Pass to $f_j\rightharpoonup f$ and, by Rellich compactness, to outputs
converging strongly in $L^2(B)$ and weakly in $H_0^1(B)$.
The preceding form-space limit puts their limit in $V_0$.
Testing the variational equation against $V_0$ identifies it as $R_0f$.
Compactness of $R_0$ also gives $R_0f_j\to R_0f$, a contradiction.
The same containing-box embedding proves compact resolvents, and the
Poincar\'e inequality in $B$ proves positivity.
\end{proof}

This argument keeps every old window fixed. In dimension one a positive
connector joins adjacent intervals for every $\delta>0$, and the point
removal argument fails; this is why $d\ge2$ is required.

Suppose $S$ is a finite set of old eigenvalues and the new cube has no
eigenvalue in $S$. Each $\lambda\in S$ is then isolated in $A_{D_0}$ with
exactly its old full eigenspace, extended by zero.
Choose disjoint isolating intervals $I_\lambda$ with endpoints outside
$\sigma(A_{D_0})$. Lemma~\ref{lem:window} implies
\begin{equation}\label{eq:clusters}
 E_{D_\delta}(I_\lambda)
   \longrightarrow E_{O_n}(\{\lambda\})\oplus0
 \quad\hbox{in operator norm}.
\end{equation}
For example, apply fixed Riesz contours around the isolated nonzero
eigenvalues $(1+\lambda)^{-1}$ of the bounded resolvents.
The resolvent identity gives convergence of the contour integrals.
Projections at norm distance below one have equal ranks: either
projection is injective on the other's range. A multiple eigenvalue may
split inside its interval; its entire cluster is retained.

\section{Inductive protection of an exhaustive family}

Set $a_1=x_1=1$. Once $O_n$ is fixed, the next cube eigenvalues are
\[
 \frac{\pi^2(m_1^2+\cdots+m_d^2)}{4a^2},
 \qquad m_i\in\N,\quad \N=\{1,2,\ldots\}.
\]
For any finite positive set $S$ and $a\in(a_n+1,a_n+2)$, an equality with
$\lambda\in S$ bounds every $m_i$. Only finitely many widths are forbidden.
We may therefore choose $a_{n+1}$ in this interval avoiding $S$, then
set $x_{n+1}=x_n+a_n+a_{n+1}$. This changes no earlier parameter.

When cube $\omega_j=Q_{a_j}(x_j)$ is chosen, fix its orthonormal sine
basis $(f_{j,\ell})_{\ell\ge1}$, extended by zero as necessary.
Let
\[
 \varepsilon_n=2^{-n-2},\qquad \eta_n=2^{-n},\qquad
 F_n=\{f_{j,\ell}:j,\ell\le n\}.
\]
We construct finite-rank orthogonal projections $P_n^k$, $1\le k\le n$,
each a sum of complete eigenspace projections of $A_{O_n}$, such that
\begin{align}
 P_n^1&\le\cdots\le P_n^n,\label{eq:nesting}\\
 \norm{P_{n+1}^k-P_n^k}&<\varepsilon_n,\qquad
 \rank P_{n+1}^k=\rank P_n^k \quad(k\le n),\label{eq:transport}\\
 \norm{(I-P_n^n)f}&<\eta_n\quad(f\in F_n).\label{eq:birth}
\end{align}
In \eqref{eq:transport} the old projections are extended by zero into the
new finite-stage $L^2$ space.

For $n=1$, a sufficiently high finite spectral cutoff gives
\eqref{eq:birth}. At stage $n$, let $S_n$ be the finite set of eigenvalues
occurring in $P_n^n$. Choose the next cube nonresonant with $S_n$.
For each $\lambda\in S_n$ take an interval as in \eqref{eq:clusters}.
If $S_n^k\subset S_n$ defines $P_n^k$, transport it by
\[
 P_{n+1}^k=\sum_{\lambda\in S_n^k}
                  E_{D_{\delta_n}}(I_\lambda).
\]
There are finitely many constraints. Choose one positive $\delta_n$ so
small that \eqref{eq:transport} holds for every $k\le n$ and
\[
 \delta_n<\min(a_n/2,2^{-n}).
\]
Set $O_{n+1}=D_{\delta_n}$. Disjoint intervals and complete clusters
preserve nesting, complete-eigenspace status and rank.
Finally choose a spectral cutoff $P_{n+1}^{n+1}$ above all transported
clusters and high enough to satisfy \eqref{eq:birth} for $F_{n+1}$.
Compact-resolvent completeness permits this finite choice.
An older projection need not remain a low-energy cutoff after transport.
No simple spectrum or uniform spectral-gap estimate is assumed.

\section{The final domain and its resolvent}

Let $\Omega=\bigcup_n O_n$. It is open and connected and contains balls
of radius $a_n\to\infty$. Apart from countably many null interface
patches it is the disjoint union of the cubes. Hence
\[
 H=L^2(\Omega)=\bigoplus_{j\ge1}L^2(\omega_j),
\]
and all $f_{j,\ell}$ together form an orthonormal basis of $H$.
Every earlier projection and resolvent now acts by zero extension on $H$.
The preceding norms and ranks are unchanged by these canonical embeddings.

Put $R=(A_\Omega+1)^{-1}$ and
$R_n=i_n(A_{O_n}+1)^{-1}r_n$, where $r_n$ restricts and $i_n$ extends
by zero. Then $R_n\to R$ strongly, with $\norm{R_n}\le1$.
Indeed the zero-extended $V_n=H_0^1(O_n)$ are increasing closed subspaces
of $V=H_0^1(\Omega)$ for
\[
 b(u,v)=\int_\Omega(\nabla u\cdot\overline{\nabla v}+u\bar v).
\]
Their union is dense: compact support in the increasing open cover lies
in one finite stage, and $C_c^\infty(\Omega)$ is form-dense.
For $f\in H$, $R_nf$ is the $b$-orthogonal projection of $Rf$ onto $V_n$.
It therefore converges in form norm and in $H$.
Norm convergence of this infinite exhaustion is not required.
The extended $R_n$ have kernels on the future cube summands; only $R$ is
used as a genuine resolvent on all of $H$.

\section{Limiting eigenvectors and completeness}

For fixed $k$, \eqref{eq:transport} gives a norm limit $P^k$ and
\begin{equation}\label{eq:tail}
 \norm{P^k-P_n^k}\le\sum_{j=n}^\infty\varepsilon_j
       =2^{-n-1}\quad(n\ge k).
\end{equation}
It is an orthogonal projection of the same finite rank, by norm
continuity and the distance-below-one criterion.
Nesting passes to the limits.
Since $R_nP_n^k=P_n^kR_n$, strong convergence of the uniformly bounded
$R_n$ and norm convergence of $P_n^k$ give
$RP^k=P^kR$.
Thus $\Ran P^k$ is a finite-dimensional reducing subspace for $R$.
Positivity and injectivity of $R$ make every eigenvalue $\mu$ on that
space strictly positive. If $Rv=\mu v$, then
$v=\mu^{-1}Rv\in\dom A_\Omega$ and
$A_\Omega v=(\mu^{-1}-1)v$.
Consequently $\Ran P^k\subset\Hpp(A_\Omega)$, the closed span of its
eigenvectors. Injectivity prevents a nonzero protected limit space from
escaping to infinite energy.

Fix $f=f_{j,\ell}$. For every $n\ge\max(j,\ell)$,
\eqref{eq:birth} and \eqref{eq:tail}, applied to the trajectory born at $n$,
give
\[
 \norm{(I-P^n)f}
 \le\norm{(I-P_n^n)f}+\norm{P_n^n-P^n}
 <2^{-n}+2^{-n-1}\longrightarrow0.
\]
Every $P^nf$ lies in the closed space $\Hpp(A_\Omega)$.
It follows that every cube-basis probe lies there and
$\Hpp(A_\Omega)=H$.
This proves absence of both continuous spectral subspaces; dense
eigenvalues alone would not prove that assertion.
Equivalently, $P^k\to I$ strongly. Diagonalizing $R$ on $\Ran P^1$ and
the finite-dimensional orthogonal differences
$\Ran(P^{k+1}-P^k)$ gives a complete orthonormal eigenbasis.

\section{Spectral support and the zero endpoint}

Nonnegativity gives $\sigma(A_\Omega)\subset[0,\infty)$.
For $\lambda\ge0$, choose $\xi\in\R^d$ with $|\xi|^2=\lambda$ and
$\chi\in C_c^\infty(B_1(0))$ of $L^2$ norm one. With cube center $c_n$
and $r_n=a_n/2$, put
\[
 u_n(x)=r_n^{-d/2}\chi((x-c_n)/r_n)e^{i\xi\cdot x}.
\]
These unit functions belong to $C_c^\infty(\Omega)\subset\dom A_\Omega$.
The product rule gives
\[
 \norm{(A_\Omega-\lambda)u_n}_2
 \le r_n^{-2}\norm{\Delta\chi}_2
       +2|\xi|r_n^{-1}\norm{\nabla\chi}_2\longrightarrow0.
\]
Their supports lie in disjoint cubes, so they are orthonormal and weakly
null. Weyl's criterion puts every $\lambda\ge0$ in the essential spectrum.
Completeness identifies the spectral set with the closure of the
eigenvalues. Finally, a zero eigenfunction has zero weak gradient and is
constant on connected $\Omega$. Its infinite volume excludes a nonzero
$L^2$ constant. This proves all assertions of Theorem~\ref{thm:main}.

\section*{Verification, priority and availability}

The accompanying source and verification archive preserves the original
proof and dated audits, independent geometric and spectral
reconstructions, bounded primary-source priority searches, their
adversarial challenge, exact finite diagnostics and file-integrity
checks. Those computations do not formally verify the infinite-dimensional
theorem or compute admissible apertures. An explicitly corrected
source-reading error in one audit is preserved with its superseding
addendum; it did not change the theorem.

The priority searches located no earlier full matching result.
They do not certify historical firstness or worldwide current openness.
Some older full bodies, a restricted finite-dimensional gap-transport
note, unindexed work and nonpublic material remained outside the
comparison. The dated question and substantive antecedents are credited
above; the claim is the complete construction proved here.
The research record is
\href{https://github.com/AlecKriebel/Math/pull/765}{repository PR~765}.
This note and its authored supporting material are released under
\href{https://creativecommons.org/licenses/by/4.0/}{CC BY 4.0}.

\begin{thebibliography}{9}
\bibitem{KL}
D. Krej\v{c}\'i\v{r}\'ik and V. Lotoreichik,
\emph{Quasi-conical domains with embedded eigenvalues},
Bull. Lond. Math. Soc. \textbf{56} (2024), 2969--2981.
\href{https://doi.org/10.1112/blms.13113}{doi:10.1112/blms.13113}.
Inspected author version:
\href{https://arxiv.org/abs/2205.08172v2}{arXiv:2205.08172v2}
(20 May 2024).
\bibitem{OWR}
V. Lotoreichik, joint work with D. Krej\v{c}\'i\v{r}\'ik,
\emph{Quasi-conical domains with embedded eigenvalues}, in
Geometric Spectral Theory, Oberwolfach Rep. \textbf{20} (2023),
report 36/2023, 2077--2078.
\href{https://doi.org/10.4171/OWR/2023/36}{doi:10.4171/OWR/2023/36}.
\bibitem{K26}
D. Krej\v{c}\'i\v{r}\'ik,
\emph{Spectral geometry: old questions and new answers},
\href{https://arxiv.org/abs/2609.28602v1}{arXiv:2609.28602v1}
(23 September 2026), Open Problem 1, PDF p.~16.
\bibitem{HSS}
R. Hempel, L. A. Seco and B. Simon,
\emph{The essential spectrum of Neumann Laplacians on some bounded
singular domains}, J. Funct. Anal. \textbf{102} (1991), 448--483.
\href{https://doi.org/10.1016/0022-1236(91)90130-W}{doi:10.1016/0022-1236(91)90130-W}.
\end{thebibliography}
\end{document}
